%===============================================================================
% APPENDIX A: Exploratory Receding-Horizon Investigation
%===============================================================================
\section{Exploratory Receding-Horizon Investigation}
\label{app:receding_horizon}

This appendix sketches one exploratory integration of RECTOR's rule evaluation and lexicographic selection layer with a pretrained \textbf{M2I DenseTNT}~\cite{Sun_2022} receding-horizon planner and serves as supplementary implementation detail for longer-horizon experiments.

\subsection{Design}

Rather than training a new trajectory generator, this integration uses three frozen, pretrained M2I components: \textbf{Relation V2V} (classifies each agent pair as influencer/reactor or independent), \textbf{DenseTNT Marginal} (generates 6~marginal trajectory modes per agent over 80~steps at 10\,Hz using VectorNet), and \textbf{Conditional V2V} (conditions reactor prediction on the influencer trajectory for interactive pair reasoning).  At each 2-second planning tick the pipeline encodes the current scene state (11-step history) via VectorNet, generates $K_{\text{plan}}{=}6$ diverse trajectory modes per agent via DenseTNT NMS, then combines marginal and conditional modes to produce 16~ego candidate trajectories per planning cycle via the CandidateGenerator wrapper.  Here, $K_{\text{plan}}{=}6$ denotes the raw DenseTNT marginal modes per agent, while the CandidateGenerator expands them to $K_{\text{ego}}{=}16$ ego trajectories through marginal and reactor-conditioned combinations.  The main RECTOR evaluation uses $K{=}6$ CVAE candidates without this conditional expansion stage.  The pipeline then evaluates all 28~rules against each mode using \texttt{KinematicRuleEvaluator}, selects the best rule-compliant mode via RECTOR's tolerance-based lexicographic ordering, and advances 20~frames (2.0\,s) before replanning with updated observations.  The rule evaluation layer (28 rules, 4 tiers, lexicographic selection) is retained unchanged.

\subsection{Planning Loop}

The full planning loop (\texttt{scripts/lib/planning\_loop.py}, 1,259~lines, 15/15 unit tests passing) integrates candidate generation, reactor selection, M2I conditional prediction, and tiered selection into a single \texttt{plan\_tick()} call with four sub-components: a \textbf{CandidateGenerator} (16~ego candidates per cycle), a \textbf{ReactorSelector} ($K{=}3$ most relevant reactors by distance, corridor overlap, and closing speed), a \textbf{CandidateScorer} (composite score using CVaR collision risk, comfort, goal progress, lane deviation, and hysteresis, with lexicographic selection as a final override), and a \textbf{KinematicRuleEvaluator} (all 28~rules evaluated via exact kinematic checks rather than learned proxies).

The pipeline was validated on 40 \texttt{validation\_interactive} scenarios from WOMD~v1.3.0 (minADE ranging from 0.40 to 3.36\,m across scenarios).  A planning tick completes in approximately 250\,ms on CPU (excluding M2I inference); target latency with GPU-batched M2I is $<$100\,ms per 2-second cycle, making real-time deployment feasible on GPU-equipped hardware.
